\begin{table}[!htbp]
\centering
\caption{Cost of delivering USD 12 million to the poorest quintile under alternative policies}
\label{tab:alternative_policies}
\begin{tabular}{lcc}
\toprule
Policy & Cost to deliver & Share of foregone revenue \\
\midrule
Perfect targeting (Q1 only) & USD 12 million & 8\% \\
Universal transfer (no targeting) & USD 60 million & 40\% \\
Universal transfer (50\% leakage) & USD 120 million & 80\% \\
Child grant and nutrition program & USD 150 million & 100\% \\
VAT exemption (actual) & USD 150 million & 100\% \\
\bottomrule
\end{tabular}
\vspace{0.3em}
\begin{minipage}{0.86\linewidth}
\footnotesize \textit{Notes:} This table compares the fiscal cost of alternative policies that would deliver the same monetary benefit to the poorest quintile as the VAT exemption reform. The benchmark assumes that the actual VAT exemption delivers approximately USD 12 million to the poorest quintile, or about 8 cents of each dollar of foregone revenue. The alternative rows report stylized transfer schemes under different targeting assumptions and implied leakage.
\end{minipage}
\end{table}
